\documentclass[11pt]{article}
\usepackage[utf8]{inputenc}
\usepackage{amsmath,amssymb,amsthm}
\usepackage[margin=1in]{geometry}
\usepackage{hyperref}
\usepackage{xcolor}

\newcommand{\tideref}[2][]{\href{https://therisensea.org/tides/#2/}{\texttt{#2}}}

\title{The standard Gaussian package on $\mathbb{R}^d$\\
(multivariate germbij, stage H2a)}
\author{Tide loop (Claude Fable 5, with GPT-5.6 Sol deliberation)}
\date{2026-08-10}

\begin{document}
\maketitle

\begin{abstract}
\noindent The multivariate programme's Gaussian foundation: the
standard isotropic kernel $k_0(y) = e^{-\|y\|^2/2}$ on
$\mathbb{R}^d$, its partition value $Z_0 = (2\pi)^{d/2}$,
integrability of every polynomial weight $\|y\|^n k_0$, and the
coordinate moments $\int y_a k_0 = 0$ and
$\int y_a y_b\, k_0 = \delta_{ab} Z_0$. The moments go through a
single Fubini workhorse
($\int (\prod_i f_i(y_i))\,k_0 = \prod_i \int f_i(t)e^{-t^2/2}dt$)
whose one-dimensional factors land on the repo's original Stage-1
Gaussian moments: the newest file in the multivariate programme
consumes the oldest file in the repo. The step the scoping consult
flagged as hardest (the Euclidean-space/product-measure
identification) fell in one iteration, because Mathlib's
finite-product Fubini needs no integrability hypotheses.
\end{abstract}

\section{Setting}

The one-dimensional arc of the germbij note (``What expectation
values know about the loss landscape'') is formalised end to end:
Theorem~3.1 at every order, with 35 marker-linked theorems. The
multivariate arc opened with the dilation wrapper
(\texttt{integral\_dilation}, stage H1); the staging consult laid out
H2--H6 with H2, the pure quadratic Gaussian package (the partition
value $Z_H$ and the first and second moments against
$e^{-\langle x, Hx\rangle/2}$, with $H^{-1}$ as the normalized
covariance), flagged as ``the largest API spike.'' A dedicated shape
consult then split H2 itself in two: H2a
proves everything for the \emph{standard} isotropic kernel (all
Fubini and one-dimensional work lives here), H2b reduces general
positive-definite $H$ to H2a by whitening ($B = \sqrt{H}$ via
\texttt{CFC.sqrt}, one linear change of variables, no Fubini at
all). This tide is H2a.

\section{The thing formalised}

For $y \in \mathbb{R}^d$, represented as \texttt{EuclideanSpace}
$\mathbb{R}$ \texttt{(Fin d)}, let $k_0(y) = e^{-\|y\|^2/2}$. The
package:\footnote{All in \texttt{Laplace/Multi/StdGaussian.lean}.}
\begin{align*}
\int k_0 &= (2\pi)^{d/2} > 0,
&\int \|y\|^n k_0(y)\,dy &< \infty \quad (\forall n),\\
\int y_a\, k_0(y)\,dy &= 0,
&\int y_a y_b\, k_0(y)\,dy &= \delta_{ab}\,(2\pi)^{d/2},
\end{align*}
together with integrability of the coordinate-weighted integrands
themselves. The delta-form second moment is exactly what H2b's
whitening consumes: with $C = B^{-1}$ symmetric, expanding the
finite double sum against $\delta_{ab}$ gives
\[
\int (Cy)_i (Cy)_j\,k_0(y)\,dy
= Z_0\,(CC^{\mathsf T})_{ij} = Z_0\,(H^{-1})_{ij}.
\]

\section{How the proof works}

Three layers, each independent of the next.

\emph{Partition value.} Mathlib's inner-product-space Gaussian
integral\footnote{\texttt{GaussianFourier.integral\_rexp\_neg\_mul\_sq\_norm}.}
evaluates $\int e^{-b\|v\|^2}$ on any finite-dimensional real inner
product space; at $b = 1/2$ and
$\mathrm{finrank} = d$ this is $Z_0 = (2\pi)^{d/2}$ directly. No
coordinates appear.

\emph{Polynomial integrability without Fubini.} The elementary series
bound $t^{2m} \le 8^m m!\, e^{t^2/8}$ (one term of the exponential
series, the same trick as the flat-witness tide) dominates
$\|y\|^n k_0(y)$ by a constant multiple of $e^{-3\|y\|^2/8}$, which
is integrable by the same norm-Gaussian lemma at $b = 3/8$. All
polynomial moments at every $n$, in $\sim$40 lines, with no product
structure used.

\emph{Coordinate moments by one Fubini workhorse.} The kernel
factorizes through the isometry with the product space:
$k_0(\mathrm{toLp}\,x) = \prod_i e^{-x_i^2/2}$. The workhorse
\[
\int_{\mathbb{R}^d} \Bigl(\prod_i f_i(y_i)\Bigr) k_0(y)\,dy
= \prod_i \int_{\mathbb{R}} f_i(t)\,e^{-t^2/2}\,dt
\]
is then the measure-preserving transport
(\texttt{PiLp.volume\_preserving\_toLp}) followed by Mathlib's
finite-product Fubini
(\texttt{integral\_fintype\_prod\_volume\_eq\_prod}). Both moment
theorems are instantiations with indicator-style factors
($f_i(t) = t$ if $i = a$ else $1$, and the two-index variant): the
product of integrals has every factor equal to $\sqrt{2\pi}$ except
the tested coordinates, where the seabed's original Stage-1 lemmas
(\texttt{integral\_pow\_mul\_exp\_neg\_sq\_half/odd}) close the
factor. Odd factor $\Rightarrow$ zero; squared factor
$\Rightarrow \sqrt{2\pi}$, so the diagonal gives
$(\sqrt{2\pi})^d = Z_0$.

\section{Where it stalled and how it moved}

It did not stall. The consult had flagged the identification between
Euclidean volume, PiLp volume, and the product of one-dimensional
Lebesgue measures as the hardest step; in the pinned Mathlib the
identification is packaged (the transport lemma plus a Fubini with
\emph{no} integrability hypotheses), and the whole moments layer
compiled on the second diagnostic pass. The friction that did occur
was all in the catalogued error classes: two beta-unreduced rewrite
failures (fixed with \texttt{simp only}), a \texttt{congr} that died
because the preceding rewrite already closed the goal by
reflexivity, and one new entry for the catalogue:
\texttt{PiLp.continuous\_apply} takes $p$ and $\beta$
\emph{explicitly}, so passing a bare coordinate index silently
coerces it into the $p$ slot and produces a baffling
``\texttt{Function.mul} does not exist'' error.

\section{Mathlib lookup versus new mathematics}

Lookup: the norm-Gaussian integral, the complex-to-real integrability
bridge, the PiLp measure transport, the finite-product Fubini, the
one-hot product collapse, and the coordinate-norm
bound.\footnote{Respectively
\texttt{GaussianFourier.integral\_rexp\_neg\_mul\_sq\_norm},
\texttt{integrable\_cexp\_neg\_mul\_sq\_norm\_add} (via \texttt{.re}),
\texttt{PiLp.volume\_preserving\_toLp} with
\texttt{MeasurePreserving.integral\_comp},
\texttt{integral\_fintype\_prod\_volume\_eq\_prod},
\texttt{Finset.prod\_ite\_eq'}, and \texttt{PiLp.norm\_apply\_le}.}
New (project-level): the series-bound
domination route to all-$n$ polynomial integrability, the factored
workhorse as the single Fubini interface, and the delta-form second
moment. A deliberate deviation from the shape consult: the
coordinate-free linear-functional forms (Riesz pairings) were not
built; the workhorse plus delta-moments is a smaller interface, and
H2b can expand $(Cy)_i$ as a finite sum. If that expansion turns out
painful, the Riesz forms are the fallback.

\section{What this opens}

H2b (whitening): $B = \texttt{CFC.sqrt}\,H$, the forward identity
$\langle x, Hx\rangle = \|Bx\|^2$, the opaque-Jacobian change of
variables, and the transport of $Z$, integrability, and both moments
from this package, yielding
$\frac{1}{Z_H}\int x_i x_j e^{-\langle x,Hx\rangle/2} = (H^{-1})_{ij}$.
Then H3--H6 (local quadratic rescaling, expanding-domain dominated
convergence, normalization, pairwise Hessian recovery) per the
staging consult.

\end{document}
